\section{Introduction}\label{sec:introduction}

The entanglement entropy of a pure quantum state measures how much information
is shared by a region and its complement. Its elementary upper bound is
proportional to the number of sites in the region. An area law replaces this
volume by the number of interactions crossing the region's boundary. We prove
such a bound for the unique gapped ground state of a general finite-range
Hamiltonian in two dimensions.

\subsection{The setting and the theorem}

Let $\Lambda\subset\Z^2$ be finite. Its edges join pairs at Euclidean distance
one, so that $\Lambda$ carries the induced square-lattice graph. Write
$d_\Lambda$ for graph distance, with value $+\infty$ between different
components. At each site put a copy of $\C^q$, and write
$\mathcal H_\Lambda=\bigotimes_{x\in\Lambda}\C^q$.
An operator supported on $X\subseteq\Lambda$ acts as the identity on the other
tensor factors. An admissible Hamiltonian has the form
\begin{equation}\label{eq:hamiltonian}
 H=\sum_{\substack{\varnothing\ne X\subseteq\Lambda\\
                  \diam_{d_\Lambda}X\le R}}h_X,
 \qquad h_X=h_X^*,\qquad \norm{h_X}\le J.
\end{equation}
There is one summand for each support $X$; a summand may be zero.
Here $q\ge1$ and $R\ge0$ are fixed integers and $J>0$ is fixed.
This indexing convention bounds the sum of interaction norms at every site
in terms of $R,J$ alone.

For $A\subseteq\Lambda$ define the edge boundary
\[
 \partial_\Lambda A=
 \{\{x,y\}:x\in A,\ y\in\Lambda\setminus A,\ |x-y|_1=1\}.
\]
If $\Omega\in\mathcal H_\Lambda$ is a unit vector, write
$\rho_A=\Tr_{\Lambda\setminus A}\ket{\Omega}\bra{\Omega}$ and
$S_\Omega(A)=-\Tr\rho_A\log\rho_A$, with $0\log0=0$.
All logarithms are natural.

\begin{theorem}[Area law]\label{thm:area}
For every $q,R,J$ as above and every $\Delta>0$, there is a finite constant
$C=C(q,R,J,\Delta)$ with the following property. Suppose that $\Lambda$ is a
finite induced square-lattice domain and that the Hamiltonian
\eqref{eq:hamiltonian} has a unique normalized ground vector $\Omega$, up to
phase, with ground energy $E_0$, and satisfies
\[
 H-E_0\id\ \ge\ \Delta\bigl(\id-\ket{\Omega}\bra{\Omega}\bigr).
\]
Then for every $A\subseteq\Lambda$,
\begin{equation}\label{eq:area-law}
 S_\Omega(A)\le C\,|\partial_\Lambda A|.
\end{equation}
The same constant applies to all such domains, sets, and Hamiltonians.
\end{theorem}

The distance in \eqref{eq:hamiltonian} is the distance within the induced
graph. This distinction matters for domains with holes or disconnected
components. The spectral assumption concerns the full Hamiltonian; no gap
assumption is imposed on Hamiltonians restricted to smaller regions.

For completeness, the theorem also gives the usual vertex-boundary
formulations. Define
\[
 \partial_{\rm in}A=\{x\in A:\exists y\notin A,
                \ \{x,y\}\in\partial_\Lambda A\},
 \qquad
 Z_A=\bigcup_{e\in\partial_\Lambda A}e.
\]

\begin{corollary}\label{cor:rectangles}
The conclusion of Theorem~\ref{thm:area} holds for every subset of a finite
rectangular square lattice with bounded on-site and nearest-neighbor
interactions and a unique ground state separated by a fixed global gap.
For every domain covered by that theorem it also implies
\[
 S_\Omega(A)\le4C|\partial_{\rm in}A|,
 \qquad S_\Omega(A)\le2C|Z_A|.
\]
\end{corollary}
\begin{proof}
On-site and nearest-neighbor interactions have graph range at most one.
Each inner boundary vertex is incident with between one and four crossing
edges, so $|\partial_{\rm in}A|\le|\partial_\Lambda A|
\le4|\partial_{\rm in}A|$. Counting both endpoints gives
$|Z_A|\le2|\partial_\Lambda A|\le4|Z_A|$.
These inequalities include empty boundaries.
\end{proof}

\subsection{History and related work}
The relation between entanglement and boundary area arose in quantum-field
calculations motivated by black-hole entropy.  Bombelli, Koul, Lee, and
Sorkin related entropy from tracing inaccessible quantum-field degrees of
freedom to black-hole area~\cite{BombelliEtAl1986}; Srednicki subsequently found area scaling
numerically for a spherical region in the vacuum of a regulated free
massless scalar field~\cite{Srednicki1993}.  These calculations motivated
the entanglement area-law question through regulated field models. For
finite-dimensional lattice spins, the gap-to-area conjecture asks whether
a unique ground state has an entropy bound proportional to boundary area
when local dimension, interaction range and strength are bounded and the
spectral gap has a size-independent positive lower
bound~\cite{Hastings2007,AradLandauVazirani2012}.
Theorem~\ref{thm:area} gives an affirmative answer in the stated
two-dimensional induced-domain setting.

In one dimension, Hastings proved a length-independent entropy bound for
the unique ground state of a finite-range Hamiltonian with uniformly
bounded interactions, fixed local dimension, and a positive spectral
gap~\cite{Hastings2007}.  Arad, Landau, and Vazirani improved the dependence
on the inverse gap for frustration-free chains using approximate ground-state
projectors~\cite{AradLandauVazirani2012}.  These are operators that preserve
the ground vector while contracting its orthogonal complement, with
controlled growth of Schmidt rank.  Arad, Kitaev, Landau, and Vazirani
extended this approach to general, possibly frustrated
chains~\cite{AradKitaevLandauVazirani2013}.  A complementary result of Brand\~ao and Horodecki
derives an area law for pure one-dimensional states from exponential
decay of correlations between arbitrary separated
regions~\cite{BrandaoHorodecki2015}.  This region-wise hypothesis is
stronger than decay of single-site correlation functions.  More generally,
the exponential clustering theorem relates a spectral gap to decay of
connected correlations of separated
observables~\cite{NachtergaeleSims2006}; an entropy estimate requires
control of the reduced state's full spectrum.
Anshu, Harrow, and Soleimanifar obtained boundary-dependent bounds on
smooth entanglement spread, which measures the width of the significant
part of that spectrum, including square-root boundary scaling on lattices
up to logarithmic factors~\cite{AnshuHarrowSoleimanifar2022}.  Such a width
estimate leaves the location of the spectrum, and hence its entropy,
to be controlled.

Strict area-law bounds in higher dimensions are also known for
restricted model classes.  Cramer, Eisert, Plenio, and Drei\ss ig proved
such a bound for arbitrary regions in cubic harmonic lattices with a
finite-range positive potential matrix whose condition number is
uniformly bounded~\cite{CramerEisertPlenioDreissig2006}.  These Gaussian
models have infinite-dimensional oscillator sites and provide a
comparison with the finite-dimensional spin setting considered here.

Other higher-dimensional results make additional spectral or thermodynamic
assumptions.  Masanes obtained bounds with logarithmic corrections from correlation
decay and bounds on the number of low-energy states of restricted
Hamiltonians~\cite{Masanes2009}; Brand\~ao and Cramer obtained a logarithmically
corrected bound for cubic regions under a low-temperature heat-capacity
hypothesis~\cite{BrandaoCramer2015}.  Cho proved strict area scaling for
balls along boundary-growth families satisfying uniform gaps and an
additional compatibility condition on successive reduced ground
states~\cite{Cho2014}.  These conditions supply information beyond the gap
of the single full Hamiltonian considered here.

For two-dimensional spin systems, Anshu, Arad, and Gosset proved a
subvolume bound for rectangular regions in a rectangular lattice with
plaquette projector interactions, a unique frustration-free ground
state, and a uniform local gap~\cite{AnshuAradGosset2022Subvolume}.
At fixed local dimension and local gap, the bound is
\(O(b^{5/3}\log^{7/3}(2+b))\), where \(b\) is the rectangle's boundary
size.  Their subsequent result gives
\[
  S\le n^{\,1+O((\log n)^{-1/5})}
\]
across a vertical cut of length \(n\), again for locally gapped,
frustration-free plaquette projector models with a unique ground
state~\cite{AnshuAradGosset2022Area}.  The local-gap assumption concerns
restrictions of the Hamiltonian; it is stronger than a gap for the full
system alone.  The latter entropy bound is near-linear, and its theorem
concerns vertical cuts.  The setting considered here allows frustrated
interactions, assumes a gap only for the full Hamiltonian, and includes
arbitrary regions in arbitrary finite induced square-lattice domains.
Arad, Firanko, and Jain extend approximate-projector methods to
maximally mixed ground states and bound their mutual information, with a
two-dimensional application under frustration-free and local-gap
assumptions~\cite{AradFirankoJain2026}.

The proof below combines locality estimates with ways of expressing entropy
through operator norms. For locality, we use the Lieb--Robinson
framework~\cite{LiebRobinson1972,NachtergaeleOgataSims2006} and Kitaev's
centered spectral filter~\cite[Appendix~D.1.2]{Kitaev2006Anyons}.
Swingle and McGreevy discuss positive constraints obtained by squaring
filtered terms, for which retaining a gap is a separate
issue~\cite[Section~8.4]{SwingleMcGreevy2016}. Here taking absolute
values gives the operator domination that retains the full-system gap.
Later, products of positive contractions derived from these constraints
preserve the ground state and suppress excitations, as projection products
do in detectability-lemma arguments~\cite{AnshuAradVidick2016}. Their localization uses propagation
estimates for unital channels, related to dissipative Lieb--Robinson
bounds~\cite{BarthelKliesch2012}; the coupled random-product estimate
needed here is proved in Section~\ref{sec:amplification}.

For the entropy comparison, bounds on local entanglement production are a
precedent for controlling logarithmic commutators independently of
spectator dimensions~\cite{VanAcoleyenMarienVerstraete2013}.
Section~\ref{sec:initial} proves the quadratic real-part estimate needed
for our marginal tails. Strong subadditivity and conditional-entropy
continuity supply the entropy inequalities used throughout
\cite{LiebRuskai1973,AlickiFannes2004,Winter2016}.
We encode marginal spectra on symmetric tensor powers using
Schur--Weyl duality~\cite{EtingofEtAl2011}, whose connection with spectrum
estimation is developed by Keyl and Werner~\cite{KeylWerner2001}.
The passage back to one-copy vectors uses finite quantum de~Finetti
measures~\cite{ChristandlEtAl2007,LewinNamRougerie2015} and the product
calculus of symmetrized finite observables~\cite{RaggioWerner1989}.
Weighted operator geometric means~\cite{PuszWoronowicz1975,KuboAndo1980}
combine the changing partitions, while the modular comparison starts
from the resolvent-compression identity of Carlen and
Vershynina~\cite[Lemma~2.1]{CarlenVershynina2019}.
The detailed construction uses these ingredients to compare entropy gain
and excitation energy in the same family of states; the next overview
explains how that comparison leads to the area bound.


\subsection{A tensor-network consequence}

A projected entangled-pair state (PEPS) is specified by one tensor at each
lattice site, with a physical index for that site's state and one virtual
index for each incident edge. Contracting the paired virtual indices gives
a many-body vector. The dimension of an edge index is its bond dimension.
The approximation problem asks how large these dimensions must be to
approximate the ground vector on the original lattice.

The companion \cite[Theorem~1.1]{OpenAI2026PolynomialPEPS} proves a
polynomial bound for the original open $L\times L$ square grid, $L\ge2$,
with on-site and nearest-neighbor interactions. For fixed local dimension
$q\ge2$, term-norm bound $J>0$, and global gap $\Delta>0$, every unique
ground vector in this class has a PEPS approximation with maximum bond
dimension at most $CL^c$ and global vector error at most $L^{-1}$ after
normalization and phase choice. The constants depend only on $q,J,\Delta$.
This is an existence result; it does not assert an efficient tensor
construction or contraction algorithm.

The approximation uses more than the entropy inequality. This paper also
provides positive operators that annihilate the ground vector and have
rapidly decaying spatial tails (Proposition~\ref{prop:positive}),
quantitative information bounds across collars
(Propositions~\ref{prop:collar} and \ref{prop:amplification}), and
two families of geometrically separated regions
(Proposition~\ref{prop:two-families}), with the region shapes and clearances
specified in those results. The companion combines these estimates and the
geometric ordering with additional compression and routing arguments;
its PEPS result is a consequence of that further analysis.

\subsection{How the proof is organized}

Fix the set $A$ whose entropy is to be bounded, and let $Z_A$ be its
crossing-edge endpoints. For disjoint physical regions $B,C$, write
\[
 I_\Omega(B:C)=S_\Omega(B)+S_\Omega(C)-S_\Omega(B\cup C)
\]
for their mutual information. Our analytic goal is to make this quantity
small when $B$ lies inside $A$ and $C$ is beyond a suitable collar,
even when $C$ includes all of $\Lambda\setminus A$.
Here is why this suffices. Suppose that, after removing a set $D$ of
$O(|\partial_\Lambda A|)$ sites, we can partition $A$ into two ordered
families of regions. If each region shares little information with the
complement of $A$ together with the earlier regions of its own family,
two applications of the entropy chain rule cancel all individual region
entropies. Purity then bounds $S_\Omega(A)$ by $|D|\log q$ and the sum
of the information errors. Section~\ref{sec:geometry} constructs regions
with the required ordering and separation.

The analytic estimates concern intersections of $A$ with ambient boxes or
finite unions of rectangles and triangles. Their clearance is from $Z_A$;
they may meet holes or the edge of the physical domain. Hamiltonian
locality and propagation always use neighborhoods in the induced graph.
The argument has five stages.
\begin{enumerate}[leftmargin=*,label=\arabic*.]
\item \emph{Obtain a first exponent below volume.}
Section~\ref{sec:initial} uses the global gap to bound fluctuations of the
negative logarithms of positive marginal eigenvalues. It then optimizes
products of positive operators on nested rectangles. Two bounds on the
optimized norm, one using the energy estimate and global gap, give a
conditional-entropy discount, and hence
\[
 S_\Omega(A\cap Q)\le Cr^{1+e_0},\qquad e_0<1,
\]
for a box $Q$ of side at most $r$ sufficiently far from $Z_A$.

\item \emph{Control information and energy in the same interpolation.}
Section~\ref{sec:quasilocal} replaces the Hamiltonian by a sum of positive
operators that annihilate $\Omega$ separately, retain a global gap, and
have rapidly decaying graph-local tails.
Sections~\ref{sec:conditional}--\ref{sec:transport} associate positive
operators to partitions of the full system and interpolate between them
as a small subsystem changes parts. These operators act on symmetric
tensor powers, whose vectors are fixed by every permutation of the copies.
Applying an inverse power of the interpolating operator to a vector built
from ground-state copies and an auxiliary entangled pair defines the
filtered vector whose norm we study. As the copy number grows, logarithmic
norms per copy can be compared with combinations of marginal entropies.
The decrease of a suitable logarithmic squared norm is bounded below by averaged conditional mutual
information, up to a small error. The same transported states bound the
energy cost. This common control permits a comparison of information
gained with excitation energy spent.

\item \emph{Compare the norms and scan a collar.}
Section~\ref{sec:comparators} gives two bounds on the negative logarithm
of the filtered vector's norm. Auxiliary dimensions are chosen from a
high-probability part of the target's Schmidt spectrum. Projecting the
target and an auxiliary system onto an entangled pair gives the upper
bound. The lower bound uses the global gap: low physical energy forces
most copies into the ground state.
Section~\ref{sec:scanner} moves the partitions through a collar according
to a finite randomized schedule. Each interaction is rarely split by a
moving boundary, and the schedule supplies an interpolation with sufficiently
low energy. The two norm bounds first improve the box exponent repeatedly;
their sharper form then gives
\[
 I_\Omega(A\cap T:\Lambda\setminus(A\cap T_L))
       \le Cn^{1-\varepsilon}.
\]
Here $T$ is one of the finite unions of controlled rectangles and triangles
specified in Definition~\ref{scanner:template}, $\varepsilon>0$ is fixed,
and $T_L$ is the ambient sup-norm enlargement by
$L=\lfloor n^{1-\varepsilon}\rfloor$. These targets and their allowed
enlargements have $O(n)$ boundary edges. The entire finite physical and
auxiliary system and schedule are fixed before the copy number tends to infinity.

\item \emph{Amplify the information bound.}
Section~\ref{sec:amplification} uses typical spectral subspaces to
construct a candidate operator for restoring a target moved to an auxiliary
reference. Its norm is controlled by the collar information. A localized
product of contractions derived from the positive constraints suppresses
the excited component, and polar decomposition yields a unitary comparison near the target. This unitary
leaves both the reference and the remote physical system untouched;
comparison with a state in which they are independent bounds their
original mutual information. Enlarging the collar by a sublinear radius
makes the bound smaller than any prescribed inverse power of $n$.

\item \emph{Arrange the regions for entropy cancellation.}
Section~\ref{sec:geometry} builds the two ordered families described above.
The collar around each region avoids the complement of $A$ and all earlier
regions of its own family. Narrow strips and recursive repairs at point
contacts produce these separations for arbitrary cuts while leaving only
$O(|\partial_\Lambda A|)$ sites uncovered. The information errors are
summable over the regions, so the two entropy chains prove
\eqref{eq:area-law}.
\end{enumerate}
